\section{Coefficient conventions, logical status, and scope}
\begin{readerroadmap}
This section sets out the conventions used throughout the chapter.  It explains when a statement concerns algebra, when topology enters the argument, which results are taken as classical input, and how the stable record IDs connect this chapter to later work.  Keeping these settings distinct makes the hypotheses of each result easier to follow.
\end{readerroadmap}
\subsection{Ambient regimes and theorem records}

\begin{definition}[Coefficient and module regimes]
\label{def:coefficient-regime}
All modules are left modules unless stated otherwise.  The algebraic regime consists of a unital ring $A$ and a unital left module $M$ over $A$.  When linear dimensions, tensor products, quivers, or operator blocks are used, $A$ is a unital algebra over $\bbk$, where $\bbk$ is a field, and $M$ is a vector space over $\bbk$.  The topological regime consists of a Hausdorff topological vector space over a complete valued field, together with a unital algebra of continuous linear operators.  Every infinite sum, completion, quotient topology, and closure below is interpreted only in the regime explicitly declared in its statement.
\end{definition}
\RecordID{PSC-I-D-COEFFICIENT-REGIME}
\paragraph{Conventions for empty families.} For finite families of nonnegative integers we use $\max\varnothing=0$. The height of an empty finite partially ordered set is $0$, as is the height of a nonempty poset with no strict comparable pair. Thus the height in the theorem for the support preorder counts strict steps.

\begin{definition}[Canonical theorem record]
\label{def:theorem-record}
A canonical theorem record is the tuple
\[
(\mathrm{ID},\mathrm{status},\mathrm{hypotheses},\mathrm{statement},
\mathrm{dependencies},\mathrm{proof},\mathrm{exports},\mathrm{boundary}).
\]
A later chapter cites the stable record ID.  A domain specialization may add domain hypotheses, but it does not repeat a generic proof from this chapter.
\end{definition}
\RecordID{PSC-I-D-THEOREM-RECORD}

\begin{boundary}[Symbolic scope]
\label{bnd:symbolic-scope}
This chapter proves the algebraic and topological statements printed below.  It does not assert historical novelty, independent external certification, automatic discovery of useful packet systems, nonlinear obstruction theory, multilevel refinement, arithmetic congruences, quantum gravity, automorphy, or local or global Langlands reciprocity.
\end{boundary}
\RecordID{PSC-I-B-SYMBOLIC-SCOPE}

\section{Typed operator grammars, realizations, and seed systems}
\begin{readerroadmap}
An operator grammar describes which compositions are allowed and which relations they satisfy.  A realization gives this syntax a concrete action on vector spaces.  Its realized algebra consists of the operators that actually occur in that action.  We can then ask which vectors serve as seeds.
\end{readerroadmap}
\subsection{Typed syntax and module realizations}

\begin{definition}[Typed operator grammar]
\label{def:typed-grammar}
Let $Q=(Q_0,Q_1,s,t)$ be a quiver with finite vertex set.  A typed operator grammar over $\bbk$ is a pair
\[
\mathfrak G=(Q,I),
\]
where $I\triangleleft \bbk Q$ is an ideal on both sides generated by linear combinations over $\bbk$ of paths having common source and common target.  Write
\[
\Lambda_{\mathfrak G}:=\bbk Q/I.
\]
A realization of $\mathfrak G$ is a family $(X_v)_{v\in Q_0}$ and linear maps
\[
X_a:X_{s(a)}\longrightarrow X_{t(a)}
\qquad(a\in Q_1)
\]
that annihilate every relation in $I$.  A morphism $f:X\to Y$ of realizations is a family of linear maps $f_v:X_v\to Y_v$ satisfying
\[
f_{t(a)}X_a=Y_af_{s(a)}
\qquad(a\in Q_1).
\]
\end{definition}
\RecordID{PSC-I-D-TYPED-GRAMMAR}

\begin{definition}[Realization on a single space]
\label{def:collapsed-realization}
For a typed realization $X$, define its realization on a single space by
\[
|X|:=\bigoplus_{v\in Q_0}X_v.
\]
Every vertex idempotent acts as a coordinate projector on $|X|$, and every path acts as a block operator from its source space to its target space.
\end{definition}
\RecordID{PSC-I-D-COLLAPSED-REALIZATION}

\begin{theorem}[Typed realization and module equivalence]
\label{thm:typed-module-equivalence}
The category of realizations of $\mathfrak G$ is equivalent to the category of unital left modules over $\Lambda_{\mathfrak G}$.  Explicitly, a realization $X$ determines the collapsed space $|X|$ of \cref{def:collapsed-realization}
with the vertex idempotent $e_v$ acting as the coordinate projector onto $X_v$ and each arrow acting in the block between its declared source and target.  Conversely, a unital module $M$ over $\Lambda_{\mathfrak G}$ determines the realization
\[
M_v=e_vM,
\qquad
M_a(x)=ax.
\]
The two constructions are inverse up to canonical natural isomorphism.
\end{theorem}
\RecordID{PSC-I-T-TYPED-MODULE-EQUIVALENCE}
\Depends{PSC-I-D-TYPED-GRAMMAR.}

\begin{proof}
Because $Q_0$ is finite, $1=\sum_{v\in Q_0}e_v$ in $\bbk Q$, the $e_v$ are pairwise orthogonal, and the same identities hold in $\Lambda_{\mathfrak G}$.  Given a realization, define the action of a path by composition of its arrow maps and extend linearly.  Every element of $I$ acts as zero, so the action factors through $\Lambda_{\mathfrak G}$.  The vertex idempotents give
\[
|X|=\bigoplus_v e_v|X|=\bigoplus_vX_v.
\]
Conversely, if $M$ is unital, then
\[
M=1M=\sum_ve_vM.
\]
Orthogonality makes the sum direct.  For an arrow $a:v\to w$, the identities $a=e_wae_v$ imply
\[
a(e_vM)\subseteq e_wM.
\]
The relations in $I$ act as zero because $M$ is a module over $\Lambda_{\mathfrak G}$.  Applying both constructions returns the original vertex spaces and arrow maps, or the original module decomposition, through the evident identity maps.  These identifications commute with morphisms, giving an equivalence of categories.
\end{proof}

\begin{interpretation}
The equivalence does not claim that a quiver presentation is unique.  It says that once a grammar is fixed, working with typed arrows and working with a module over the quotient path algebra carry the same information.  The typed description is often easier for constructing an operator system.  The module description is often easier for proving generation and quotient theorems.
\end{interpretation}

\begin{microexample}[title={A grammar with two vertices}]
Take two vertices and one arrow \(a:1\to2\).  A realization is a pair of spaces and a linear map
\[
X_1\xrightarrow{\,X_a\,}X_2.
\]
After collapse, \(|X|=X_1\oplus X_2\), the vertex idempotents are the two coordinate projectors, and \(a\) is a single block outside the diagonal.  This elementary case already contains the viewpoint of packets and transfer blocks used later.
\end{microexample}


\subsection{The realized action and faithful reduction}
\begin{definition}[Realized operator algebra]
\label{def:realized-algebra}
For a realization $X$, the realized operator algebra is
\[
\A_X:=\im\!\left(\Lambda_{\mathfrak G}\longrightarrow\End_{\bbk}(|X|)\right).
\]
It records the actual operators in the realization and may be a proper quotient of the abstract grammar algebra.
\end{definition}
\RecordID{PSC-I-D-REALIZED-ALGEBRA}

\begin{proposition}[Action factorization and hull invariance]
\label{prop:action-factorization}
Let $S\subseteq |X|$.  The generated submodule under the grammar algebra equals the generated submodule under the realized operator algebra:
\[
\Lambda_{\mathfrak G}S=\A_XS.
\]
Consequently, every statement about seed rank depends only on the realized action, not on redundant relations in a chosen presentation of the grammar.
\end{proposition}
\RecordID{PSC-I-P-ACTION-FACTORIZATION}
\Depends{PSC-I-D-REALIZED-ALGEBRA.}

\begin{proof}
The action map $\rho:\Lambda_{\mathfrak G}\twoheadrightarrow\A_X$ is surjective by definition and satisfies $\lambda s=\rho(\lambda)s$.  Hence every grammar orbit vector lies in $\A_XS$.  The reverse inclusion holds because every element of $\A_X$ is $\rho(\lambda)$ for some $\lambda$.
\end{proof}

\begin{keyidea}
Replacing the abstract grammar algebra by its image \(\A_X\) removes syntax that acts as zero.  This is a faithful reduction of the action.  It does not change any orbit, hull, or computation of seed rank.
\end{keyidea}


\subsection{Seeds, addressed hulls, and natural presentations}
\begin{definition}[Linear seed system]
\label{def:seed-system}
A linear seed system is a tuple
\[
\mathfrak S=(X,\A,S,\Pcal,\tau),
\]
where $X$ is a unital left module over $\A$, $\A$ is fixed independently of the target element, $S=(s_1,\ldots,s_r)$ is a finite distinguished seed family, $\Pcal$ is either empty or a declared packet system, and $\tau$ is either the algebraic convention or a specified Hausdorff topology in which the declared operator action is continuous.
\end{definition}
\RecordID{PSC-I-D-SEED-SYSTEM}

\begin{definition}[Generated hull]
\label{def:generated-hull}
For $S\subseteq X$, define
\[
\Hull_{\A}(S):=\sum_{s\in S}\A s.
\]
In the topological regime define
\[
\Hull_{\A,\tau}(S):=\cl{\Hull_{\A}(S)}^{\,\tau}.
\]
For a singleton $S=\{s\}$, write $\Hull_{\A}(s)$.
\end{definition}
\RecordID{PSC-I-D-GENERATED-HULL}

The algebraic hull uses finite sums of addressed seed vectors.  In the topological setting, taking its closure may add limits that no finite expression reaches.  Thus algebraic generation and dense generation answer different questions, even when they use the same seeds and operators.

\begin{definition}[Natural presentation]
\label{def:natural-presentation}
A seed presentation is natural relative to a declared grammar when the primitive operators and relations are fixed independently of the target, the relation system is finite or effectively enumerable, the topology is intrinsic to the realization, no primitive operator directly stores the complete target, and the presentation exposes a reusable invariant, theorem, or computation not encoded by a primitive chosen for the target.
\end{definition}
\RecordID{PSC-I-D-NATURAL-PRESENTATION}

\begin{proposition}[Algebra generation and module generation are distinct]
\label{prop:algebra-module-generation}
Every nonzero unital algebra $A$, regarded as a left module over itself, has
\[
\srank_A(A)=1,
\]
with cyclic seed $1_A$.  This conclusion is independent of the number of algebra generators required to present $A$ as an algebra.
\end{proposition}
\RecordID{PSC-I-P-ALGEBRA-MODULE-GENERATION}
\Depends{PSC-I-D-GENERATED-HULL.}

\begin{proof}
For every $a\in A$, one has $a=a1_A$.  Hence $A1_A=A$.  The number of operations needed to generate $A$ as an algebra is a different invariant and does not enter the module calculation.
\end{proof}

\begin{boundary}[Operators chosen for a target]
\label{bnd:target-dependent-operators}
For a fixed seed $s$ and target $x$, the existence of a specially defined operator $A_x$ satisfying $A_xs=x$ does not establish a natural seed calculus unless $A_x$ belongs to a grammar that is fixed independently of the target and satisfies \cref{def:natural-presentation}.
\end{boundary}
\RecordID{PSC-I-B-TARGET-DEPENDENT-OPERATORS}

\begin{warningbox}
An operator may be chosen because it sends the current seed to a desired target.  Such a choice proves only an existence statement for that one target unless the operator already belongs to the fixed algebra.  The algebra must be declared before the question of seed rank is asked.  Otherwise the acting system changes with the answer one wants to obtain.
\end{warningbox}


\begin{figure}[htbp]
\centering
\includegraphics[width=0.96\textwidth]{01_seed_operator_grammar.pdf}
\caption{Schematic organization of the seed, operator, relation, and reconstruction data. This is a guide to the presentation, not an additional formal theorem.}
\label{fig:grammar}
\end{figure}

\section{Free presentations by seeds and relations}
\begin{readerroadmap}
A seed family gives a canonical map from a finite free module.  The generated hull is exactly its image, and the relations among addressed seeds form the kernel.  This converts orbit generation into a presentation problem.
\end{readerroadmap}
\subsection{Finite free covers}
\begin{definition}[Seed evaluation map]
\label{def:seed-evaluation}
Let $A$ be a unital ring, let $M$ be a left module over $A$, and let $S=(s_1,\ldots,s_r)\in M^r$.  Define the linear map over $A$
\[
\varepsilon_S:A^{\oplus r}\longrightarrow M,
\qquad
(a_1,\ldots,a_r)\longmapsto\sum_{j=1}^{r}a_js_j.
\]
\end{definition}
\RecordID{PSC-I-D-SEED-EVALUATION}

\begin{theorem}[Finite free presentation]
\label{thm:free-presentation}
For the map in \cref{def:seed-evaluation},
\[
\im\varepsilon_S=\Hull_A(S),
\qquad
N_S:=\Ker\varepsilon_S\leq A^{\oplus r},
\]
and the induced map is an isomorphism of left modules over $A$
\[
\boxed{A^{\oplus r}/N_S\cong\Hull_A(S).}
\]
Consequently, $S$ generates $M$ if and only if $\varepsilon_S$ is surjective.
\end{theorem}
\RecordID{PSC-I-T-FREE-PRESENTATION}
\Depends{PSC-I-D-SEED-EVALUATION and the first isomorphism theorem.}
\Exports{All later free seed presentations.}

\begin{proof}
The image consists exactly of finite sums $\sum_ja_js_j$, which is $\Hull_A(S)$.  Because $\varepsilon_S$ is linear over $A$, its kernel is a left submodule.  The first isomorphism theorem gives
\[
A^{\oplus r}/N_S\cong\im\varepsilon_S=\Hull_A(S).
\]
The final equivalence is the equality $\im\varepsilon_S=M$.
\end{proof}

\begin{interpretation}
The theorem converts the geometric language of addressed seed orbits into the algebraic language of generators and relations.  A word in the operators is represented by an element of the free module.  Two addressed expressions represent the same vector exactly when their difference lies in the kernel.  The generated hull is therefore not merely spanned by the free module.  It is its canonical quotient.
\end{interpretation}

\begin{microexample}[title={A cyclic nilpotent chain}]
Let \(A=\bbk[t]\) and \(M=\bbk[t]/(t^3)\).  The residue class \(s=[1]\) is one seed and
\[
\varepsilon_s:A\to M,
\qquad
f(t)\mapsto f(t)[1]
\]
has kernel \((t^3)\).  Hence \(M\cong A/(t^3)\), and the addressed vectors \(s,ts,t^2s\) are the visible orbit chain.
\end{microexample}


\begin{corollary}[Characterization by minimal free covers]
\label{cor:minimal-free-cover}
For every finitely generated left module $M$ over $A$,
\[
\boxed{
\srank_A(M)=\min\{r\geq0:\text{there exists a surjection }A^{\oplus r}\twoheadrightarrow M\}.
}
\]
\end{corollary}
\RecordID{PSC-I-C-MINIMAL-FREE-COVER}
\Depends{PSC-I-T-FREE-PRESENTATION.}

\begin{proof}
A seed family of size $r$ gives a surjection by \cref{thm:free-presentation}.  Conversely, the images of the standard basis vectors under any surjection $A^{\oplus r}\twoheadrightarrow M$ form a generating family of size at most $r$.
\end{proof}

\begin{corollary}[Presentation with one seed]
\label{cor:one-seed-presentation}
For $s\in M$, the annihilator
\[
\operatorname{Ann}_A(s):=\{a\in A:as=0\}
\]
is a left ideal and
\[
\boxed{A/\operatorname{Ann}_A(s)\cong\Hull_A(s).}
\]
\end{corollary}
\RecordID{PSC-I-C-ONE-SEED-PRESENTATION}
\Depends{PSC-I-T-FREE-PRESENTATION.}

\begin{proof}
The evaluation map for one seed is $a\mapsto as$.  Its kernel is $\operatorname{Ann}_A(s)$, which is therefore a left ideal, and \cref{thm:free-presentation} gives the quotient isomorphism.
\end{proof}

\subsection{One seed and the distinction between module and algebra relations}
\begin{proposition}[Criterion for the annihilator to be an ideal on both sides]
\label{prop:two-sided-annihilator}
For $s\in M$, the left ideal $L=\operatorname{Ann}_A(s)$ is an ideal on both sides if and only if
\[
L\,\Hull_A(s)=0.
\]
If $s$ is cyclic, these conditions are equivalent to
\[
\operatorname{Ann}_A(s)=\operatorname{Ann}_A(M).
\]
\end{proposition}
\RecordID{PSC-I-P-ANNIHILATOR-TWO-SIDED}
\Depends{PSC-I-C-ONE-SEED-PRESENTATION.}

\begin{proof}
The left ideal $L$ is an ideal on both sides exactly when $ab\in L$ for every $a\in L$ and $b\in A$.  This is equivalent to
\[
0=(ab)s=a(bs)
\qquad(a\in L,\ b\in A).
\]
By additivity, the last identity is equivalent to $L\Hull_A(s)=0$.

If $s$ is cyclic, then $\Hull_A(s)=M$.  Hence the preceding condition is $L\subseteq\operatorname{Ann}_A(M)$.  The reverse inclusion always holds because an operator annihilating all of $M$ annihilates $s$.  Thus equality follows.  Conversely, $\operatorname{Ann}_A(M)$ is an ideal on both sides, so equality implies that $L$ is an ideal on both sides.
\end{proof}

\begin{microexample}[title={Why a seed annihilator need not be an ideal on both sides}]
Let \(A=M_2(\bbk)\) act on \(M=\bbk^2\), and let \(s=e_1\).  Then \(\operatorname{Ann}_A(s)\) consists of matrices whose first column is zero.  It is a left ideal because left multiplication preserves a zero first column.  It is not generally a right ideal.  Thus the natural quotient for one seed is a quotient of the left regular module, not automatically a quotient algebra.
\end{microexample}


\subsection{Topological completion of a presentation}
\begin{theorem}[Completion of a dense free presentation]
\label{thm:completed-presentation}
Let $X$ be a complete Hausdorff locally convex space, let $\A\subseteq\End(X)$ act continuously, and let $S=(s_1,\ldots,s_r)$.  Put
\[
E:=\A^{\oplus r}/N_S
\]
algebraically and let $j:E\to X$ be the injective map induced by $\varepsilon_S$.  Give $E$ the locally convex topology transported from $j(E)$, equivalently the topology generated by
\[
\widetilde p(e):=p(j(e))
\]
for the continuous seminorms $p$ on $X$.  Every $a\in\A$ acts continuously on $E$ in the transported topology and therefore extends to $\widehat E$.  If $\Hull_{\A}(S)$ is dense in $X$, then $j$ extends uniquely to a topological isomorphism of modules over $\A$
\[
\boxed{\widehat E\cong X.}
\]
\end{theorem}
\RecordID{PSC-I-T-COMPLETED-PRESENTATION}
\Depends{PSC-I-T-FREE-PRESENTATION and uniqueness of Hausdorff completion.}

\begin{proof}
By \cref{thm:free-presentation}, $j$ is an algebraic isomorphism of modules over $\A$ from $E$ onto $\Hull_{\A}(S)$.  The transported seminorms make it a topological isomorphism onto that subspace.  For $a\in\A$ and a continuous seminorm $p$ on $X$, continuity of $a:X\to X$ supplies continuous seminorms controlling $p(aj(e))$, so the induced action on $E$ is continuous and extends uniquely to $\widehat E$.  The subspace is dense by hypothesis.  A Cauchy net in $E$ is mapped to a Cauchy net in $X$, which converges because $X$ is complete.  This defines a continuous extension $\widehat j:\widehat E\to X$.  For any $x\in X$, choose a net $(e_\lambda)$ in $E$ with $j(e_\lambda)\to x$.  The transported topology makes $(e_\lambda)$ Cauchy, so it converges to some $\xi\in\widehat E$, and continuity gives $\widehat j(\xi)=x$.  Thus $\widehat j$ is surjective.  If $\widehat j(\xi)=0$, represent $\xi$ by a Cauchy net $(e_\lambda)$.  Then $j(e_\lambda)\to0$, so every defining seminorm $\widetilde p(e_\lambda)$ tends to zero.  Thus $e_\lambda\to0$ in the completion and $\xi=0$.  Hence $\widehat j$ is bijective.  Its inverse is continuous by the universal property of completion, or directly because the topology on $E$ was transported from $X$.
\end{proof}

The topology in this completion theorem comes from the realized image.  It is not an independently chosen topology on the free algebra.  This choice ensures that the completion recovers the given space, while leaving separate the question of whether another natural topology on the grammar produces the same completion.

\begin{sectionrecap}
A seed family always gives a free presentation of its hull.  The kernel stores all linear relations among addressed seed words.  One seed gives a cyclic module, but only special annihilators permit an algebra quotient.  Dense generation is handled by completing the presented image rather than by silently replacing algebraic span with closure.
\end{sectionrecap}

\section{Seed rank and functorial calculus}



\begin{readerroadmap}
For a finitely generated module, seed rank is the smallest number of seeds needed to generate it.  We first examine how this number behaves when the acting algebra changes or the module is replaced by a quotient.  We then give bounds for exact sequences, direct sums, and tensor products.  Quotients with simple actions provide useful witnesses for lower bounds.
\end{readerroadmap}
\subsection{Definition, cyclicity, and monotonicity}

\begin{definition}[Seed rank]
\label{def:seed-rank}
For a unital left module $M$ over $A$, define
\[
\srank_A(M):=\min\{|S|:AS=M\}.
\]
For a topological module $(X,\tau)$, define
\[
\srank_{\A,\tau}(X):=\min\{|S|:\cl{\A S}^{\,\tau}=X\}.
\]
Set $\srank_A(0)=0$.  When the value is finite, it is the minimal number of module generators.
\end{definition}
\RecordID{PSC-I-D-SEED-RANK}

\begin{proposition}[Cyclic characterization]
\label{prop:cyclic-characterization}
For $M\neq0$,
\[
\srank_A(M)=1
\quad\Longleftrightarrow\quad
M\text{ contains a cyclic vector}.
\]
\end{proposition}
\RecordID{PSC-I-P-CYCLIC-CHARACTERIZATION}
\Depends{PSC-I-D-SEED-RANK.}

\begin{proof}
Both conditions say that there is $m\in M$ with $Am=M$.
\end{proof}

\begin{proposition}[Monotonicity under enlargement of the acting algebra]
\label{prop:acting-algebra-monotonicity}
Let $B\subseteq A\subseteq\End_{\bbk}(X)$ be unital operator algebras with the same identity.  Then
\[
\boxed{\srank_A(X)\leq\srank_B(X).}
\]
The same statement holds for topological seed rank when both algebras act continuously.
\end{proposition}
\RecordID{PSC-I-P-ACTING-ALGEBRA-MONOTONICITY}
\Depends{PSC-I-D-SEED-RANK.}

\begin{proof}
For every $S\subseteq X$, one has $BS\subseteq AS$.  Hence every generating family over $B$ is also a generating family over $A$.  Taking closures gives the topological statement.
\end{proof}

\begin{interpretation}
Enlarging the acting algebra can only create more orbit directions, so it can only lower the number of seeds required.  This is why every statement about seed rank must name the acting algebra.  The same vector space may be cyclic for a large algebra and require many seeds for a smaller one.
\end{interpretation}


\subsection{Quotients, isomorphisms, and exact sequences}
\begin{theorem}[Quotient monotonicity]
\label{thm:quotient-monotonicity}
Let $\varphi:A\to B$ be a unital ring homomorphism and let
\[
\Phi:M\twoheadrightarrow N
\]
be a surjective additive map satisfying
\[
\Phi(am)=\varphi(a)\Phi(m)
\qquad(a\in A,\ m\in M).
\]
Then
\[
\boxed{\srank_B(N)\leq\srank_A(M).}
\]
If $M$ and $N$ are topological modules, $\Phi$ is continuous and surjective, and the actions are continuous, then
\[
\srank_{B,\tau_N}(N)\leq\srank_{A,\tau_M}(M).
\]
\end{theorem}
\RecordID{PSC-I-T-QUOTIENT-MONOTONICITY}
\Depends{PSC-I-D-GENERATED-HULL and PSC-I-D-SEED-RANK.}
\Exports{All quotient lower bounds and realization inequalities.}

\begin{proof}
Let $S$ generate $M$.  For every $n\in N$, choose $m\in M$ with $\Phi(m)=n$ and write $m=\sum_ia_is_i$.  Then
\[
n=\sum_i\varphi(a_i)\Phi(s_i),
\]
so $\Phi(S)$ generates $N$.

For the topological statement, let $S$ generate $M$ densely.  Given $n=\Phi(m)$, choose a net $m_\lambda\in AS$ converging to $m$.  Continuity gives
\[
\Phi(m_\lambda)\longrightarrow n,
\]
and every $\Phi(m_\lambda)$ lies in $B\Phi(S)$.  Thus $B\Phi(S)$ is dense in $N$.
\end{proof}

\begin{microexample}[title={A quotient can strictly lower seed rank}]
Let \(A=\bbk\) act by scalars on \(M=\bbk^3\).  Then \(\srank_A(M)=3\).  Quotienting by a subspace of dimension one gives a scalar module of dimension two, whose seed rank is \(2\).  Quotient monotonicity is therefore often strict.
\end{microexample}


\begin{corollary}[Isomorphism invariance]
\label{cor:isomorphism-invariance}
An intertwining module isomorphism preserves algebraic seed rank.  A topological intertwining isomorphism preserves topological seed rank.
\end{corollary}
\RecordID{PSC-I-C-ISOMORPHISM-INVARIANCE}
\Depends{PSC-I-T-QUOTIENT-MONOTONICITY.}

\begin{proof}
Apply \cref{thm:quotient-monotonicity} to the isomorphism and its inverse.
\end{proof}

\begin{theorem}[Short exact sequence bounds]
\label{thm:exact-sequence-bounds}
Let
\[
0\longrightarrow L\xrightarrow{\iota}M\xrightarrow{\pi}N\longrightarrow0
\]
be a short exact sequence of left modules over $A$.  Then
\[
\boxed{
\srank_A(N)\leq\srank_A(M)
\leq\srank_A(L)+\srank_A(N).
}
\]
\end{theorem}
\RecordID{PSC-I-T-EXACT-SEQUENCE-BOUNDS}
\Depends{PSC-I-T-QUOTIENT-MONOTONICITY.}

\begin{proof}
The lower inequality is quotient monotonicity. For the upper inequality, choose generating families $(\ell_i)_{i\in I}$ of $L$ and $(n_j)_{j\in J}$ of $N$ with the least possible cardinalities. Choose a lift $m_j\in M$ of each $n_j$. For any $m\in M$, a finite subset $F\subseteq J$ and coefficients $a_j\in A$ give
\[
\pi(m)=\sum_{j\in F}a_jn_j.
\]
The vector $m-\sum_{j\in F}a_jm_j$ belongs to $\iota(L)$, so it is a finite linear combination of the $\iota(\ell_i)$. The union of the lifted family and the image of the generating family of $L$ therefore generates $M$. Its cardinality is at most $|I|+|J|$, giving the upper bound. This argument uses finite sums for each vector even when either generating family is infinite.
\end{proof}

\begin{hypothesisguide}
For a short exact sequence, a generating family for \(M\) always maps to a generating family for the quotient.  It need not generate the submodule \(L\).  For the upper bound, choose generators of \(L\) separately and combine them with lifts of quotient generators.
\end{hypothesisguide}


\subsection{Direct sums under shared and independent actions}
\begin{proposition}[Direct sum bounds for a shared action]
\label{prop:shared-direct-sum}
If one algebra $A$ acts diagonally on $M\oplus N$, then
\[
\boxed{
\max\{\srank_A(M),\srank_A(N)\}
\leq\srank_A(M\oplus N)
\leq\srank_A(M)+\srank_A(N).
}
\]
\end{proposition}
\RecordID{PSC-I-P-SHARED-DIRECT-SUM}
\Depends{PSC-I-T-QUOTIENT-MONOTONICITY.}

\begin{proof}
The coordinate projections are intertwining quotients and give the lower bound.  For the upper bound, embed minimal generating families of $M$ as $(m_i,0)$ and of $N$ as $(0,n_j)$. Their union generates by finite sums and has cardinality at most the sum of the two seed ranks.
\end{proof}

\begin{theorem}[Direct sum formula for independent actions]
\label{thm:independent-direct-sum}
Let $M$ be a left module over $A$ and $N$ a left module over $B$.  Under the action of the product algebra
\[
(a,b)(m,n)=(am,bn),
\]
one has
\[
\boxed{
\srank_{A\times B}(M\oplus N)
=
\max\{\srank_A(M),\srank_B(N)\}.
}
\]
\end{theorem}
\RecordID{PSC-I-T-INDEPENDENT-DIRECT-SUM}
\Depends{PSC-I-T-QUOTIENT-MONOTONICITY.}

\begin{proof}
The coordinate projections give the lower bound. Let $r$ be the maximum on the right, regarded as a cardinal, and choose an index set $I$ of cardinality $r$. Choose minimal generating families for $M$ and $N$, index them by subsets of $I$, and fill any unused positions with zeros. Write the resulting families as $(m_i)_{i\in I}$ and $(n_i)_{i\in I}$. For any $m\in M$ and $n\in N$, their finite expressions involve a common finite subset $F\subseteq I$ after zero coefficients have been added. Thus
\[
m=\sum_{i\in F}a_im_i,
\qquad
n=\sum_{i\in F}b_in_i,
\]
and
\[
(m,n)=\sum_{i\in F}(a_i,b_i)(m_i,n_i).
\]
The family of paired vectors generates and has cardinality at most $r$. The empty family handles the case in which both modules are zero.
\end{proof}

\subsection{Tensor generation and witnesses for lower bounds}
\begin{proposition}[Tensor product bound]
\label{prop:tensor-bound}
Let $A$ and $B$ be unital algebras over $\bbk$, let $M$ be a left module over $A$, and let $N$ be a left module over $B$.  Under the natural $A\otimes_{\bbk}B$ action,
\[
\boxed{
\srank_{A\otimes B}(M\otimes_{\bbk}N)
\leq
\srank_A(M)\srank_B(N).
}
\]
\end{proposition}
\RecordID{PSC-I-P-TENSOR-BOUND}
\Depends{PSC-I-D-SEED-RANK.}

\begin{proof}
Choose minimal generating families $(m_i)_{i\in I}$ for $M$ and $(n_j)_{j\in J}$ for $N$. Each vector in either module is a finite linear combination of its generators. Bilinearity therefore expresses every simple tensor as a finite sum of terms
\[
(am_i)\otimes(bn_j)=(a\otimes b)(m_i\otimes n_j).
\]
Simple tensors span the algebraic tensor product, so the family $(m_i\otimes n_j)_{(i,j)\in I\times J}$ generates it. This family has cardinality at most $|I||J|$, which is the claimed bound. The proof also covers an empty generating family, when the corresponding module and tensor product are zero.
\end{proof}

\begin{interpretation}
The tensor product bound is constructive.  Tensor the chosen seed families and act independently in the two factors.  Equality can fail because a single tensor seed may contain several useful components, or because the tensor product algebra acquires additional relations.  Later quantum applications should therefore cite the bound, not assume multiplicativity.
\end{interpretation}


\begin{boundary}[Tensor rank is not generally multiplicative]
\label{bnd:tensor-nonequality}
The inequality in \cref{prop:tensor-bound} can be strict.  Equality requires additional hypotheses on the acting algebras and their module multiplicities and is not exported as a universal law.
\end{boundary}
\RecordID{PSC-I-B-TENSOR-NONEQUALITY}



\begin{proposition}[Monotonicity under extension of scalars]
\label{prop:scalar-extension}
Let \(K/\bbk\) be a field extension, let \(A\) be a unital algebra over \(\bbk\), and let \(M\) be a finitely generated left module over \(A\).  Set
\[
A_K:=K\otimes_{\bbk}A,
\qquad
M_K:=K\otimes_{\bbk}M.
\]
Then
\[
\boxed{
\srank_{A_K}(M_K)\leq\srank_A(M).
}
\]
Every generating family \(S\subseteq M\) extends to the generating family \(1\otimes S\subseteq M_K\).
\end{proposition}
\RecordID{PSC-I-P-SCALAR-EXTENSION-MONOTONICITY}
\Depends{PSC-I-D-SEED-RANK.}
\Exports{Extension of the coefficient field in realizations from arithmetic, quantum theory, and representation theory.}

\begin{proof}
Let \(S=\{s_1,\ldots,s_r\}\) generate \(M\).  Every pure tensor \(\lambda\otimes m\in M_K\) has an expression
\[
m=\sum_{i=1}^r a_is_i,
\]
so
\[
\lambda\otimes m
=
\sum_{i=1}^r(\lambda\otimes a_i)(1\otimes s_i).
\]
Pure tensors span \(M_K\), hence \(1\otimes S\) generates \(M_K\) over \(A_K\).  Taking the minimum over \(S\) proves the inequality.
\end{proof}

\begin{hypothesisguide}
The proposition gives an inequality in one direction.  A smaller generating family found after extending scalars need not descend to the original field.  Descent requires additional structure.  The chapter therefore exports monotonicity, not an unconditional equality.
\end{hypothesisguide}


\subsection{Quotient witnesses for lower bounds}
\begin{definition}[Witness for seed rank]
\label{def:seed-rank-witness}
A surjective intertwiner $\Phi:M\twoheadrightarrow N$ is a witness for seed rank of order $r$ when $\srank(N)=r$.  It certifies $\srank(M)\geq r$.
\end{definition}
\RecordID{PSC-I-D-SEED-RANK-WITNESS}

\begin{proposition}[Scalar quotient witness]
\label{prop:scalar-quotient}
Let $M$ be a module over $A$ with coefficient field $\bbk$.  Suppose $M$ has a quotient $N$ of dimension $r$ on which $A$ acts through a character
\[
\chi:A\longrightarrow\bbk.
\]
Then
\[
\boxed{\srank_A(M)\geq r.}
\]
\end{proposition}
\RecordID{PSC-I-P-SCALAR-QUOTIENT-WITNESS}
\Depends{PSC-I-T-QUOTIENT-MONOTONICITY.}

\begin{proof}
For every $n\in N$, the cyclic submodule $An$ is contained in the scalar span $\bbk n$.  Therefore $q$ seed images span a subspace of dimension at most $q$.  Generating the quotient of dimension $r$ requires at least $r$ seeds.  Quotient monotonicity transfers the lower bound to $M$.
\end{proof}

\begin{sectionrecap}
Seed rank is functorial under quotients and isomorphisms, subadditive across extensions, given by the maximum for independent product actions, and bounded above under tensor products.  Lower bounds come from quotients on which the action simplifies.  Scalar extension cannot increase seed rank, but descent of a reduced seed family is a separate problem.
\end{sectionrecap}

\section{Generator number modulo the Jacobson radical}



\begin{readerroadmap}
For a finitely generated module, Nakayama's lemma allows us to count generators after passing to the quotient by the Jacobson radical.  When the algebra has finite dimension, this quotient is semisimple and its multiplicities give an explicit answer.  We begin with the radical reduction and then calculate those multiplicities.
\end{readerroadmap}
\subsection{Reduction to the module top}

\begin{classicalinput}[Nakayama lemma]
\label{ci:nakayama}
Let $A$ be a unital ring, let $J\subseteq\Jac(A)$, and let $M$ be a finitely generated left module over $A$.  If $M=N+JM$ for a submodule $N\leq M$, then $M=N$.
\end{classicalinput}
\RecordID{PSC-I-CI-NAKAYAMA}

\begin{theorem}[Generator number is detected modulo the Jacobson radical]
\label{thm:module-top}
Let $A$ be a unital ring, let $J=\Jac(A)$, and let $M$ be a finitely generated left module over $A$.  A family $S\subseteq M$ generates $M$ if and only if its image generates $M/JM$ as a module over $A/J$.  When $A/J$ is semisimple Artinian, this quotient is the usual semisimple top.  No semisimplicity of $A/J$ is required for the general equality.  Consequently,
\[
\boxed{
\srank_A(M)=\srank_{A/J}(M/JM).
}
\]
\end{theorem}
\RecordID{PSC-I-T-MODULE-TOP}
\Depends{PSC-I-CI-NAKAYAMA and PSC-I-T-QUOTIENT-MONOTONICITY.}
\Exports{Generator formulas in finite dimension, quantum multiplicity calculations, and residual deformation modules.}

\begin{proof}
If $S$ generates $M$, its image generates every quotient.  Conversely, suppose the image of $S$ generates $M/JM$ and put $N=AS$.  Then
\[
M=N+JM.
\]
Nakayama's lemma gives $M=N$.  Thus the same finite families generate before and after passage to the top, and the minimal cardinalities agree.
\end{proof}

\begin{interpretation}
The theorem says that nilpotent or radical directions can lengthen orbit chains but cannot reduce the number of independent directions in the module top that must be seeded.  Once the images of the seeds generate \(M/JM\), Nakayama's lemma forces the entire module to be generated.
\end{interpretation}

\begin{microexample}[title={A local nilpotent algebra}]
Let \(A=\bbk[\varepsilon]/(\varepsilon^3)\) and \(M=A^{\oplus2}\).  The radical is \(J=(\varepsilon)\), and
\[
M/JM\cong\bbk^2.
\]
Therefore two seeds are necessary and sufficient.  The extra nilpotent layers \(\varepsilon M\) and \(\varepsilon^2M\) are generated automatically once the two top directions are seeded.
\end{microexample}


\begin{corollary}[Local algebra formula]
\label{cor:local-algebra-formula}
Let $A$ be a local ring with Jacobson radical $J$ and residue division ring $D=A/J$.  For every finitely generated left module $M$ over $A$,
\[
\boxed{
\srank_A(M)=\dim_D(M/JM),
}
\]
where the dimension is that of a left vector space over $D$.
\end{corollary}
\RecordID{PSC-I-C-LOCAL-ALGEBRA-FORMULA}
\Depends{PSC-I-T-MODULE-TOP.}

\begin{proof}
The module $M/JM$ over $A/J$ is a left vector space over $D$.  Its minimal number of generators is its dimension over $D$.
\end{proof}

\subsection{Semisimple multiplicity calculus}
\begin{classicalinput}[Wedderburn and Artin decomposition]
\label{ci:wedderburn}
Every semisimple Artinian ring is isomorphic to a finite product
\[
A\cong\prod_{i=1}^{q}M_{d_i}(D_i)
\]
of full matrix rings over division rings.
\end{classicalinput}
\RecordID{PSC-I-CI-WEDDERBURN-ARTIN}

\begin{lemma}[Multiplicities in the regular module]
\label{lem:regular-multiplicities}
Let
\[
A=\prod_{i=1}^{q}M_{d_i}(D_i)
\]
and let $S_i=D_i^{d_i}$ denote the natural simple left module for the $i$th factor, extended by zero across the other factors.  Then, as a left module over $A$,
\[
\boxed{A\cong\bigoplus_{i=1}^{q}S_i^{\oplus d_i}.}
\]
\end{lemma}
\RecordID{PSC-I-L-REGULAR-MODULE-MULTIPLICITIES}
\Depends{PSC-I-CI-WEDDERBURN-ARTIN.}

\begin{proof}
For one factor $M_d(D)$, let $e_{jj}$ be the diagonal matrix units.  The left regular module decomposes as
\[
M_d(D)=\bigoplus_{j=1}^{d}M_d(D)e_{jj}.
\]
Each minimal left ideal $M_d(D)e_{jj}$ is isomorphic to the natural column module $D^d$.  Taking the direct product of factors gives the stated decomposition.
\end{proof}

\begin{theorem}[Semisimple seed rank formula over arbitrary division rings]
\label{thm:semisimple-seed-rank}
Let
\[
A\cong\prod_{i=1}^{q}M_{d_i}(D_i)
\]
be semisimple Artinian and let
\[
M\cong\bigoplus_{i=1}^{q}S_i^{\oplus m_i}
\]
be a finitely generated left module over $A$.  Then
\[
\boxed{
\srank_A(M)
=
\max_{1\leq i\leq q}
\left\lceil\frac{m_i}{d_i}\right\rceil.
}
\]
If $A$ is any algebra of finite dimension over a field, the same formula applied to the decomposition of $M/\Jac(A)M$ over $A/\Jac(A)$ computes $\srank_A(M)$.
\end{theorem}
\RecordID{PSC-I-T-SEMISIMPLE-SEED-RANK}
\Depends{PSC-I-T-MODULE-TOP, PSC-I-CI-WEDDERBURN-ARTIN, PSC-I-L-REGULAR-MODULE-MULTIPLICITIES, and PSC-I-C-MINIMAL-FREE-COVER.}
\Exports{Every calculation of seed rank for semisimple algebras of finite dimension in the series.}

\begin{proof}
By \cref{cor:minimal-free-cover}, $M$ is generated by $r$ elements exactly when it is a quotient of $A^{\oplus r}$.  By \cref{lem:regular-multiplicities},
\[
A^{\oplus r}\cong\bigoplus_iS_i^{\oplus rd_i}.
\]
Every short exact sequence of modules over a semisimple ring splits.  Therefore a quotient isomorphic to $M$ exists if and only if
\[
m_i\leq rd_i
\qquad\text{for every }i.
\]
The smallest such $r$ is the displayed maximum.  For a nonsemisimple algebra of finite dimension, \cref{thm:module-top} reduces the calculation to the semisimple module $M/\Jac(A)M$ over $A/\Jac(A)$.
\end{proof}

\begin{interpretation}
For the factor \(M_{d_i}(D_i)\), one free generator contributes \(d_i\) copies of the corresponding simple module.  If the top of \(M\) contains \(m_i\) copies, at least \(\lceil m_i/d_i\rceil\) free generators are required.  The worst simple type controls the global answer.

This count concerns copies of a simple module, rather than the number of coordinates in a chosen basis.  A single seed can reach several copies through the matrix action.  The regular module records precisely how many copies one free generator can supply.
\end{interpretation}

\begin{microexample}[title={Multiplicity consumes seed capacity}]
For \(A=M_2(\bbk)\) acting on \((\bbk^2)^{\oplus3}\), the regular module contributes two copies of the natural simple module per generator.  One generator cannot cover all three multiplicity directions.  Two can, giving seed rank \(2\).
\end{microexample}


\Needspace{16\baselineskip}
\begin{corollary}[Algebraically closed multiplicity formula]
\label{cor:algebraically-closed-formula}
Let $\bbk$ be algebraically closed, let $A$ be a unital algebra of finite dimension over $\bbk$, put $J=\Jac(A)$, and let $M$ be a finitely generated left module over $A$.  Assume compatible decompositions indexed by an integer $q\geq0$:
\[
A/J\cong\prod_{i=1}^{q}M_{d_i}(\bbk),
\qquad
M/JM\cong\bigoplus_{i=1}^{q}(\bbk^{d_i})^{\oplus m_i},
\]
where the $d_i$ are positive integers and the $m_i$ are nonnegative integers.  The first isomorphism is an isomorphism of algebras over $\bbk$.  The second is an isomorphism of left modules over $A/J$, using the natural column module for each matrix factor and zero action from the other factors.  Then
\[
\srank_A(M)=\max_{1\leq i\leq q}\left\lceil\frac{m_i}{d_i}\right\rceil,
\]
with $\max\varnothing=0$ when $q=0$.
\end{corollary}
\RecordID{PSC-I-C-ALGEBRAICALLY-CLOSED-FORMULA}
\Depends{PSC-I-T-MODULE-TOP and PSC-I-T-SEMISIMPLE-SEED-RANK.}

\begin{proof}
By \cref{thm:module-top}, $\srank_A(M)=\srank_{A/J}(M/JM)$.  Apply \cref{thm:semisimple-seed-rank} to the displayed decompositions of $A/J$ and $M/JM$, with every $D_i=\bbk$.  This gives the stated maximum.  If $q=0$, the unital $A/J$-module $M/JM$ is zero, and \cref{thm:module-top} gives seed rank zero.
\end{proof}

\begin{corollary}[Full matrix algebra]
\label{cor:full-matrix}
For the natural module $S=D^d$ over $M_d(D)$,
\[
\srank_{M_d(D)}(S)=1.
\]
More strongly, every nonzero vector of $S$ is cyclic.
\end{corollary}
\RecordID{PSC-I-C-FULL-MATRIX}
\Depends{PSC-I-T-SEMISIMPLE-SEED-RANK.}

\begin{proof}
The formula gives $\lceil1/d\rceil=1$.  Directly, let $v\neq0$ and choose a coordinate $v_j\neq0$.  For any target $w=(w_i)$, the matrix whose $j$th column has entries $w_iv_j^{-1}$ and whose other columns vanish sends $v$ to $w$.
\end{proof}

\begin{corollary}[Multiplicity obstruction]
\label{cor:multiplicity-obstruction}
For the diagonal action of $M_d(D)$ on $S^{\oplus m}$,
\[
\boxed{
\srank_{M_d(D)}(S^{\oplus m})
=
\left\lceil\frac{m}{d}\right\rceil.
}
\]
\end{corollary}
\RecordID{PSC-I-C-MULTIPLICITY-OBSTRUCTION}
\Depends{PSC-I-T-SEMISIMPLE-SEED-RANK.}

\begin{proof}
Apply \cref{thm:semisimple-seed-rank} to the single simple type with multiplicity $m$.
\end{proof}

\begin{sectionrecap}
For finitely generated modules, the minimal number of seeds is detected on the quotient by the Jacobson radical.  In the semisimple top, each generator has a fixed multiplicity capacity \(d_i\) in simple type \(i\), and the largest capacity ratio determines the exact seed rank.
\end{sectionrecap}

\begin{figure}[htbp]
\centering
\includegraphics[width=0.96\textwidth]{04_semisimple_seed_rank.pdf}
\caption{The semisimple multiplicity formula for one simple block. A seed can expose at most $d$ independent multiplicity directions, so $m$ copies require $\lceil m/d\rceil$ seeds.}
\label{fig:semisimple}
\end{figure}

\section{Modules for a single operator over arbitrary fields}



\begin{readerroadmap}
One operator is equivalent to a module over \(\bbk[t]\).  Primary decomposition isolates the irreducible polynomial sectors, and the quotient by \(p(T)\) counts how many cyclic chains begin in each sector.
\end{readerroadmap}
\subsection{Primary components and generator counts}

\begin{classicalinput}[Primary decomposition for polynomial modules]
\label{ci:primary-decomposition}
Every module $X$ over $\bbk[t]$ with finite dimension over the coefficient field decomposes as
\[
X=\bigoplus_{p\in\Irr(\bbk[t])}X_p,
\]
where $\Irr(\bbk[t])$ denotes the monic irreducible polynomials, only finitely many components are nonzero, and every nonzero $X_p$ is a finite direct sum of modules $\bbk[t]/(p^e)$.
\end{classicalinput}
\RecordID{PSC-I-CI-PID-PRIMARY-DECOMPOSITION}

\begin{theorem}[Seed rank for a single operator]
\label{thm:single-operator}
Let $X$ have finite dimension over $\bbk$ and let $T\in\End_{\bbk}(X)$.  Regard $X$ as a module over $\bbk[t]$ by $t\cdot x=Tx$.  For every monic irreducible $p$, let $X_p$ be the primary component for $p$ and set
\[
\mu_p(T):=
\dim_{\bbk[t]/(p)}\bigl(X_p/p(T)X_p\bigr).
\]
Then
\[
\boxed{
\srank_{\bbk[T]}(X)=\max_p\mu_p(T).
}
\]
Moreover, $\mu_p(T)$ is the number of cyclic primary summands for $p$ in the decomposition into elementary divisors of $X_p$.
\end{theorem}
\RecordID{PSC-I-T-SINGLE-OPERATOR}
\Depends{\nolinkurl{PSC-I-CI-PID-PRIMARY-DECOMPOSITION},
\nolinkurl{PSC-I-C-LOCAL-ALGEBRA-FORMULA}, and
\nolinkurl{PSC-I-T-INDEPENDENT-DIRECT-SUM}.}
\Exports{Companion, recurrence, and systems of subsequences in residue classes.}

\begin{proof}
Let the minimal polynomial be
\[
m_T(t)=\prod_{p}p(t)^{e_p}.
\]
The acting algebra is
\[
\bbk[T]\cong\bbk[t]/(m_T)
\cong\prod_pR_p,
\qquad
R_p:=\bbk[t]/(p^{e_p}),
\]
by the Chinese remainder theorem.  The factor $R_p$ acts on $X_p$ and annihilates every other primary component.  Hence the formula for a direct sum with independent actions gives
\[
\srank_{\bbk[T]}(X)=\max_p\srank_{R_p}(X_p).
\]
Each $R_p$ is local with maximal ideal generated by the class of $p$ and residue division field $K_p=\bbk[t]/(p)$.  By \cref{cor:local-algebra-formula},
\[
\srank_{R_p}(X_p)
=
\dim_{K_p}X_p/p(T)X_p
=
\mu_p(T).
\]
Finally, every summand $\bbk[t]/(p^e)$ contributes one copy of $K_p$ after quotienting by $p$, so the dimension equals the number of summands in the elementary divisor decomposition.
\end{proof}

\begin{interpretation}
For one operator, the quotient \(X_p/p(T)X_p\) counts the number of independent primary chains for the irreducible polynomial \(p\).  Over an algebraic closure these become Jordan chains.  The seed rank is the largest chain count among the primary sectors, not the largest chain length.
\end{interpretation}

\begin{microexample}[title={Block count versus block size}]
A nilpotent operator with Jordan block sizes \((5,1,1)\) has seed rank \(3\), even though its largest block has length \(5\).  Three seeds are needed because there are three independent chains.  Once their initial vectors are chosen, repeated application of \(T\) generates every position in those chains.
\end{microexample}


\subsection{Invariant factors, Jordan blocks, and cyclicity}
\begin{corollary}[Invariant factor count]
\label{cor:invariant-factor-count}
If $X\neq0$ and the rational canonical decomposition is written with nonconstant invariant factors
\[
X\cong\bigoplus_{j=1}^{r}\bbk[t]/(f_j),
\qquad
f_1\mid f_2\mid\cdots\mid f_r,
\]
then
\[
\boxed{\srank_{\bbk[T]}(X)=r.}
\]
\end{corollary}
\RecordID{PSC-I-C-INVARIANT-FACTOR-COUNT}
\Depends{PSC-I-T-SINGLE-OPERATOR.}

\begin{proof}
For each irreducible $p$, the number of invariant factors divisible by $p$ equals the number of elementary divisors for $p$.  Because $f_1\mid\cdots\mid f_r$ and $f_1\neq1$ after zero summands are removed, some irreducible divisor of $f_1$ divides all $r$ invariant factors.  Hence $\max_p\mu_p(T)=r$.
\end{proof}

\begin{corollary}[Formula for geometric multiplicity]
\label{cor:geometric-multiplicity}
If $\bbk$ is algebraically closed, then
\[
\boxed{
\srank_{\bbk[T]}(X)
=
\max_{\lambda\in\operatorname{Spec}(T)}
\dim_{\bbk}\Ker(T-\lambda I).
}
\]
\end{corollary}
\RecordID{PSC-I-C-ALGEBRAICALLY-CLOSED-SINGLE-OPERATOR}
\Depends{PSC-I-T-SINGLE-OPERATOR.}

\begin{proof}
The irreducibles are $p=t-\lambda$.  On the primary component for $\lambda$, the quotient by $(T-\lambda I)X_\lambda$ has dimension equal to the number of Jordan blocks, which is also $\dim\Ker(T-\lambda I)$.
\end{proof}

\begin{corollary}[Criterion for a cyclic matrix]
\label{cor:cyclic-matrix}
For $X\ne0$ of finite dimension, the following are equivalent:
\begin{enumerate}[label=(\roman*)]
\item $T$ has a cyclic vector.
\item $\srank_{\bbk[T]}(X)=1$.
\item Every nonzero primary component for $p$ is cyclic.
\item The rational canonical form has one invariant factor.
\item The minimal polynomial equals the characteristic polynomial.
\item After scalar extension to an algebraic closure, every eigenvalue has geometric multiplicity one.
\end{enumerate}
\end{corollary}
\RecordID{PSC-I-C-CYCLIC-MATRIX}
\Depends{PSC-I-T-SINGLE-OPERATOR, PSC-I-C-INVARIANT-FACTOR-COUNT, and PSC-I-P-CYCLIC-CHARACTERIZATION.}

\begin{proof}
The equivalence of (i) and (ii) is \cref{prop:cyclic-characterization}.  By \cref{thm:single-operator}, (ii) holds exactly when every $\mu_p(T)$ belonging to a nonzero primary component equals one, which is (iii).  Relatively prime primary cyclic modules combine into one cyclic module by the Chinese remainder theorem, so (iii) is equivalent to (iv).  In rational canonical form, one invariant factor means that the unique invariant factor is both the minimal and characteristic polynomial, proving (iv)$\Leftrightarrow$(v).  Scalar extension preserves the equality of minimal and characteristic polynomials.  Over an algebraic closure, that equality is equivalent to one Jordan block for each eigenvalue, which is (vi).
\end{proof}

\begin{sectionrecap}
A single operator is cyclic exactly when each primary sector has one cyclic chain.  Over an algebraic closure this is equivalent to one Jordan block per eigenvalue and to equality of the minimal and characteristic polynomials.
\end{sectionrecap}

\begin{figure}[htbp]
\centering
\includegraphics[width=0.96\textwidth]{03_cyclic_rank_profiles.pdf}
\caption{Krylov rank profiles for a cyclic companion operator and for two equal primary blocks. The second orbit stalls before reaching the ambient dimension because one seed cannot separate the repeated blocks.}
\label{fig:krylov}
\end{figure}

